
\chapter{Inverse, implicit function theorems and submanifolds}

We will now get to the hard part of the course (it's so over). For functions of class $C^{1}(I)$ where $I \subset \mathbb{R}$ is an interval, around any point, where $f'(x_{0}) \neq 0$, we can take a neighborhood small enough around $x_{0}$ such that $f$ is monotone on that neighborhood. The inverse function theorem and proposition about the derivative of the inverse guarantee us the existence and properties of an inverse to $f$ on the small neighborhood.

The goal of this chapter is to find a similar theorem for multivariable functions. This time instead of $f'(x_{0}) \neq 0$ we will require that $\det Jf(x_{0}) \neq 0$, a very natural generalisation. We begin with a very special case.

\section{Small perturbations of the identity}

We begin by proving forms of the inverse function theorem and auxiliary lemmas for special cases which we will employ later to prove the general ones. If a function from $\mathbb{R}^{n}$ to $\mathbb{R}^{n}$ is very close to the identity, is it invertible then?

\begin{lemma}[Small Lipschitz perturbations of the identity]\label{lem:lipschitz-perturb}
	Let $U \subset \mathbb{R}^{n}$ be open and assume that $F: U \to \mathbb{R}^{m}$ is a function of the form $F(x) = x + \phi(x)$, where $\phi: U \to \mathbb{R}^{m}$ is $\lambda$-Lipschitz with $\lambda < 1$.
	\begin{enumerate}
		\item For any $x_{0} \in U \text{ s.t. } B_{r}(x_{0}) \subset U$ with some radius $r > 0$, we have
		      \begin{align*}
			      B_{(1-\lambda)r}(F(x_{0})) \subset F(B_{r}(x_{0})) \subset F(U).
		      \end{align*}
		      In particular, $F(U)$ is open.
		\item $F: U \to F(U)$ is bijective and $F^{-1}$ is $\frac{1}{1-\lambda}$-Lipschitz.
	\end{enumerate}
\end{lemma}
\begin{proof}
	To prove (1), given $y \in B_{(1-\lambda)r}(F(x_{0}))$, we will show that $\exists x \in B_{r}(x_{0})$ s.t. $F(x) = y$, which holds if and only if $x + \phi(x) = y$. So we want to find the $x$ such that $x = y - \phi(x) = T(x)$, for the map $T(x) = y - \phi(x)$. So we're looking for fixed points of $T$. But $T$ is a contraction map, as for any $x, x' \in U$
	\begin{align*}
		\left| T(x) - T(x') \right| = \left| y - \phi(x) - y + \phi(x') \right| = \left| \phi(x) - \phi(x') \right| \leq \lambda \left| x - x' \right|.
	\end{align*}
    And so we would love to apply the Banach Fixed Point Theorem \ref{thm:banach_fixed_point}. But we can't apply it exactly, as we're not sure that the elements of our sequence remain in the domain. We can however copy parts of the proof. 

    So fix $x_{0} \in U$, some radius $r > 0$ and take any $y \in B_{(1-\lambda)r}(F(x_{0}))$. Consider the sequence defined by
    \begin{align*}
        x_{k+1} = T(x_{k}) = y - \phi(x_{k}).
    \end{align*}
    This is of course a priori not well-defined, as we still need to show $x_{k} \in U$ and once we know that, we wish to show that the sequence has a limit. Let us now do that.

    Notice first that 
    \begin{align*}
        \left| x_{1} - x_{0} \right| = \left| y - \phi(x_{0})  - x_{0} \right| = \left| y - F(x_{0}) \right| < (1-\lambda)r.
    \end{align*}
    This shows $\left| x_{0} - x_{1} \right| < r$, and we can proceed to define $x_{2}$. We have 
    \begin{align*}
        \left| x_{2} - x_{1} \right| = \left| T(x_{1}) - T(x_{0}) \right| \leq \lambda\left| x_{1} - x_{0} \right| < \lambda(1 - \lambda)r
    \end{align*}
    and one checks that $x_{2} \in B_{r}(x_{0})$ as well, by triangle inequality. Doing this over and over, we arrive at
    \begin{align*}
        \left| x_{k+1} - x_{k} \right| \leq \lambda^{k} \left| x_{1} - x_{0} \right|,
    \end{align*}
    which implies
    \begin{align*}
        \left| x_{k+1} - x_{0} \right| \leq \sum_{i=0}^{k}  \left| x_{i+1} - x_{i} \right| \leq \left| x_{1} - x_{0} \right| \sum_{i=0}^{\infty}\lambda^{i} < (1-\lambda) r \sum_{i=0}^{\infty}\lambda^{i} = r
    \end{align*}
    hence $x_{k} \in B_{r}(x_{0})$ for all $k$. By the same argument as in the Banach Fixed Point Theorem, we get a point $x \in \mathbb{R}^{n}$ that satisfies $\left| x - x_{0} \right| < r$, is hence in the ball $B_{r}(x_{0})$ and that is a fixed point of the map $T$, so satisfies $\phi(x) + x = y$. 

    Now we are in position to prove (2). Assume $F(x) = y$ and $F(x') = y$ for some $x,x' \in U$. This happens if and only if
    \begin{align*}
        x + \phi(x) = x' + \phi(x') \iff \left| x-x' \right| = \left| \phi(x) - \phi(x') \right| \leq \lambda \left| x - x' \right|,
    \end{align*}
    and as $\lambda < 1$, this implies $x = x'$.

    Now we prove the Lipschitz bound of the inverse. We have $F(x) = y$ and $F(x') = y'$, so 
    \begin{align*}
        y - y' &= x - x' + \phi(x) - \phi(x') \\ 
        \implies \left| y - y' \right| &\geq \left| x - x' \right| - \left| \phi(x) - \phi(x') \right| 
        \leq \left| x-x' \right| - \lambda\left| x-x' \right| = (1-\lambda)\left| x-x' \right|.
    \end{align*}
    But this means 
    \begin{align*}
        \left| F^{-1}(y) - F^{-1}(y') \right| = \left| x-x' \right| \leq \frac{1}{1-\lambda} \left| y - y' \right|
    \end{align*}
    and $F^{-1}$ is $\frac{1}{1-\lambda}$-Lipschitz.
\end{proof}

\newpage

\begin{lemma}[Automatic differential of the inverse]\label{lem:auto-diff-inv}
    Let $U, V \subset \mathbb{R}^{n}$ be open and $f: U \to V$ and $g: V \to U$ be bijective maps such that $f(g(y)) = y$ for any $y \in V$. Assume that $f$ is differentiable at some point $x_{0} \in U$, the differential $Df_{x_{0}}: \mathbb{R}^{n} \to \mathbb{R}^{n}$ is \textit{invertible} and $g$ is $\Lambda$-Lipschitz. Then, $g$ is differentiable at the point $y_{0} = f(x_{0})$ and 
    \begin{align*}
       Dg_{y_{0}} = (Df_{x_{0}})^{-1}.
    \end{align*}
\end{lemma}
\begin{proof}
    If $f$ is differentiable at $x_{0}$, we have
    \begin{align*}
        f(x) - f(x_{0}) = Df_{x_{0}}(x-x_{0}) + \smalO (x - x_{0}),
    \end{align*}
    as $x \to x_{0}$. Plugging in $x = g(y)$ for some $y \in V$
    \begin{align*}
        y - y_{0} = f(g(y)) - f(g(y_{0})) = Df_{x_{0}}\Big(g(f(y)) - g(f(y_{0}))\Big) + \smalO( g(y) - g(y_{0}) )
    \end{align*}
    as $y \to y_{0}$. We can apply $Df_{x_{0}}^{-1}$ on both sides of this and rearrange to get
    \begin{align*}
        g(y) - g(y_{0}) = Df_{x_{0}}^{-1}(y - y_{0}) + Df_{x_{0}}^{-1} \Big( \smalO(g(y) - g(y_{0})) \Big).
    \end{align*}
    And so to prove the theorem, it remains to show that the last term is $\smalO(y - y_{0})$. But 
    \begin{align*}
        Df_{x_{0}}^{-1} \Big( \smalO ( g(y) - g(y_{0}) ) \Big) = \smalO \Big( Df_{x_{0}}^{-1}( g(y) - g(y_{0}))  \Big) = \smalO \Big(  \lVert Df_{x_{0}}^{-1} \rVert_{2} \cdot \Lambda \left| y - y_{0} \right|  \Big) = \smalO (y - y_{0}),
    \end{align*}
    using where $\lVert \cdot  \rVert_{2}$ is the Hilbert-Schmidt norm, we applied Lemma \ref{lem:hilbert-schmidt-norm} and that $g$ is $\Lambda$-Lipschitz. Also note that it's not hard to prove that $\smalO$ and linear maps commute (again, using properties of the HS norm).
\end{proof}

\begin{lemma}[Smoothness of the inverse]\label{lem:smoothness-inverse}
    Let $U \subset \mathbb{R}^{n \times n}$ be the set of invertible real $n \times n$ matrices and let $\Theta: U \to U: X \mapsto X^{-1}$ be the inverse map. Then $\Theta$ is $C^{\infty}(U, \mathbb{R}^{n \times n})$.
\end{lemma}
\begin{proof}
    Consider that the inverse of $X$ is given by the classical adjoint divided by the determinant $\det (X) \neq 0$. Since component-wise, $X^{-1}$ is hence expressed as a quotient, product, sum and composition of polynomials of components of $X$, it is indeed infinitely differentiable, using in particular that no division by zero will ever occur ($\det (X) \neq 0$).
\end{proof}



